\section{Discussion and limitations}\label{sec:discussion}

\subsection{What the decomposition does and does not establish}

The unit of analysis is a model-pair gain rather than a model, and that is what makes the
two failure modes visible. A model may estimate within-group label frequencies better
without predicting any individual case better, and be credited as though it had done the
second. Conversely it may discriminate cases better while losing enough prior fit to hide
the improvement --- Section~\ref{sec:vgvisual} contains a case where the aggregate score
and the accuracy both point the wrong way. The signed identity $\Delta T=\Delta P+\Delta R$
exposes both without a fitted reference forecaster.

Everything is conditional on the declared $\phi$. A positive $\Delta R$ means the new model
matches its probability changes to the correct case better than the old one \emph{among
cases sharing $\phi$}; it proves nothing about causal inference, compositional
understanding, or robustness. If an unrecorded shortcut varies inside a cell, its predictive
use is counted in $\Delta R$, and Proposition~\ref{prop:sensitivity} bounds that exposure
rather than removing it --- at a level that, for our own Visual Genome estimates, is not
comfortable (Appendix~\ref{app:sensitivity}). Applications should declare $\phi$ before
seeing results, justify it from how the benchmark was built, and report finer and coarser
alternatives, because coarsening adds a between-cell covariance term of either sign
(Proposition~\ref{prop:coarsen}) and only the finest identified partition nets out every
regularity expressible in $\phi$. Even that rule is weaker than it looks:
Appendix~\ref{app:partition} shows two equally fine, fully identified partitions of the same
four cases that disagree completely about whether an improvement is case-level or
transported, so the choice of $\phi$ is a modelling decision and not a formality.

Two limits deserve stating flatly. First, $\Delta R$ is not invariant to monotone
recalibration: it is a covariance between probability movements and labels, so sharpening or
softening a fixed model shifts score between the channels while adding no information about
any case, changing no ranking and no top-1 decision. Section~\ref{sec:simulation} measures
the effect at $-0.488$ in a design where both models carry identical case-level
information --- tens of times any real gain reported here --- which is why we audit
confidence-matched pairs and why released checkpoints, which routinely differ sharply in
calibration, cannot be compared on their raw outputs. A rank-based within-cell contrast
would be invariant to this, at the cost of the exact additivity that is the point of the
construction; Appendix~\ref{app:partition} says what that trade would involve. Second, a
transported gain is not evidence of anything illegitimate. Fitting the deployment label
distribution better is real skill whenever that distribution is the one that matters; it is
fragile precisely under label shift, where it can be added or removed post hoc by
re-weighting \citep{saerens2002adjusting,lipton2018detecting}, as
Section~\ref{sec:vgvisual} demonstrates and Section~\ref{sec:decision} exploits. Nor is it
purely a matter of prior fit (Proposition~\ref{prop:transported}). The decomposition reports
where a gain lives, not whether it is worth having.

\subsection{Choosing and declaring the grouping variable}

``Declare $\phi$ before looking'' is advice, not a mechanism. Our recommendation is that
the \emph{benchmark}, not the submitting author, fix it: the annotation prior is a property
of the dataset, and a maintainer-declared $\phi$ removes the incentive to report only the
most favourable candidate. That only moves the discretion, however. Maintainers are often
authors of methods evaluated on their own benchmark, and a $\phi$ that flatters the
benchmark's own baseline is no better than an author-chosen one. What we would defend is the
weaker claim that $\phi$ should be fixed and public \emph{before} a leaderboard opens, by
whoever does it, and that every pre-declared candidate be reported as our sensitivity rows
do. Guarantees holding uniformly over a class of groupings are the subject of
multicalibration \citep{hebert2018multicalibration}, and porting that style of guarantee to
pairwise attribution is a natural extension. To make reports comparable we suggest a fixed
block per audited pair: the declared $\phi$ and its rationale, the identified coverage, both
channels signed with the covariance interval, the confidence-matched row, and the
granularity rows.

\subsection{What the estimator cannot do}

Exact transport needs at least two cases per cell, so a very fine $\phi$ buys confound
control at the cost of coverage; we report the identified fraction and claim nothing for
singletons. Because cells identified at one granularity need not coincide with those at
another, cross-granularity comparisons mix Proposition~\ref{prop:coarsen}'s between-cell
term with a change of identified subpopulation --- immaterial at VG150's $99.0\%$ coverage,
but worth restricting to a common subsample when coverage is low. Approximate matching would
extend the method to continuous priors but would trade the exact identity for modelling
assumptions, and we leave that open.

The image-clustered interval is asymptotic. Theorem~\ref{thm:clt} justifies it in a regime
no heavy-tailed benchmark attains in sample: of $6{,}346$ eligible class-pair cells,
$2{,}286$ hold fewer than five cases, though those carry only $2.7\%$ of identified
relations, so most of the evidence comes from cells well above the minimum. The reported
$94.0\%$ empirical coverage against a nominal $95\%$ (Section~\ref{sec:simulation})
therefore remains the operative evidence for the interval, with the theorem supplying the
limiting justification the simulation checks against. The randomization test is
finite-group exact only under a sharper within-cell exchangeability null and treats
relations rather than images as the randomized units, so it stays a secondary diagnostic; an
image-level design preserving several relation-cell margins at once would strengthen it.

Finally, the estimator reads probability vectors, not top-1 decisions. That is a feature ---
it detects useful redistribution below the argmax --- but it requires released logits or
probabilities, and this is the real obstacle to routine use. The audit of
Section~\ref{sec:sggaudit} needed nothing but two checkpoints evaluated with their authors'
code; the method consumed the resulting arrays and returned a decomposition with intervals
in seconds. Nothing in that workflow is specific to scene graph generation. What blocks it
elsewhere is neither compute nor complexity but disclosure, since benchmarks overwhelmingly
archive ranked predictions. A leaderboard that stored per-example probabilities would make
this kind of attribution available for every pair it ranks.

\subsection{Broader applicability}\label{sec:broader}

Four aligned arrays are all the algorithm requires -- two models' probability vectors,
the labels, and the prior-cell identifiers -- and nothing in
Sections~\ref{sec:setup}--\ref{sec:inference} is specific to relation prediction.
Appendix~\ref{app:lt} puts that generality to the test on long-tailed image
classification, where taking the benchmark's coarse superclasses as $\phi$ lets the
decomposition distinguish a re-weighting schedule whose near-zero aggregate score hides
two large cancelling channels from one whose genuine improvement is part prior-refitting
and part rare-class covariance. The claim of breadth thereby acquires empirical content
instead of resting on analogy. That second domain is also where the constraint on $\phi$
shows its teeth, since the fine class label cannot serve as a cell: it would leave
every cell label-homogeneous and drive $\Delta R$ to zero by construction, producing an
artifact that reads exactly like a null result. Elsewhere the natural choices are easy
enough to name -- a predeclared question template for VQA, an answer format or prompt
family for other multimodal benchmarks -- though we have not evaluated them, and we flag
them as the obvious next domains rather than claim a coverage we have not earned.
Whatever the setting, the scientific claim has to stay precise, because what WAGER
measures is new prediction--label alignment within the declared cells and nothing more.
Its value lies not in any single $\phi$ settling what counts as understanding, but in
making every attribution an explicit and reproducible counterfactual tied to a declared
nuisance structure, one that the same unmodified construction has now answered sensibly
on two structurally different tasks.
